\documentclass[11pt]{amsart}

\usepackage[margin=1.15in]{geometry}
\usepackage{amsmath,amssymb,amsthm,mathtools}
\usepackage{microtype}
\usepackage[colorlinks=true,linkcolor=blue,citecolor=blue,urlcolor=blue]{hyperref}
\usepackage[nameinlink,capitalize]{cleveref}

\theoremstyle{plain}
\newtheorem{theorem}{Theorem}[section]
\newtheorem{proposition}[theorem]{Proposition}
\newtheorem{lemma}[theorem]{Lemma}
\newtheorem{corollary}[theorem]{Corollary}

\theoremstyle{definition}
\newtheorem{definition}[theorem]{Definition}
\newtheorem{hypothesis}[theorem]{Hypothesis}

\theoremstyle{remark}
\newtheorem{remark}[theorem]{Remark}

\numberwithin{equation}{section}

\title[Structured ancient Navier--Stokes solutions]{Structured bounded ancient solutions of the three-dimensional Navier--Stokes equations}
\author{Guillem Duran Ballester}
\date{}

\begin{document}

\begin{abstract}
We isolate three structured classes of ancient solutions of the three-dimensional
Navier--Stokes equations which are natural in the analysis of Type I singularities:
axisymmetric solutions without swirl, axisymmetric solutions with quantitative
swirl control, and solutions with uniform decay in self-similar variables.  The
axisymmetric no-swirl case is reduced to the Liouville theorem of
Koch--Nadirashvili--Seregin--Sverak by a time-shift argument which is compatible
with the Type I bound.  The other two classes are formulated as precise
Liouville hypotheses with stated admissible parameter ranges.  Assuming those
hypotheses, the corresponding nonzero Type I ancient limits cannot occur.  Thus
the paper separates an unconditional classical case from two cases whose
resolution requires genuinely new Liouville theorems.
\end{abstract}

\maketitle

\section{Introduction}

Let \(U:\mathbb R^{3}\times(-\infty,0)\to\mathbb R^{3}\) and
\(P:\mathbb R^{3}\times(-\infty,0)\to\mathbb R\) solve
\begin{equation}\label{eq:ns}
\partial_t U-\Delta U+U\cdot\nabla U+\nabla P=0,\qquad
\nabla\cdot U=0.
\end{equation}
The Type I analysis of a finite-time singularity leads, after rescaling about a
singular point, to nonzero ancient solutions satisfying
\begin{equation}\label{eq:type-i}
\sup_{t<0}\sqrt{-t}\,\|U(\cdot,t)\|_{L^\infty(\mathbb R^3)}<\infty .
\end{equation}
Here and below, a nonzero ancient limit is understood in the usual blow-up sense:
the rescalings have been normalized so that the limiting solution is not
identically zero, for instance through a positive local scale-invariant norm.
The preceding papers in this series reduce the exclusion of Type I singularities
to the exclusion of nonzero ancient limits satisfying increasingly restrictive
compactness and rigidity hypotheses.  The present paper treats additional
geometric and spatial structures that are not consequences of compactness alone.

The first structure is axisymmetry without swirl.  In this case the argument is
closed by an existing Liouville theorem: bounded ancient axisymmetric no-swirl
solutions are constant.  Since a solution satisfying \eqref{eq:type-i} is bounded
on every time interval \((-\infty,T]\), \(T<0\), this theorem applies after a
time shift.  The Type I decay as \(t\to-\infty\) then eliminates the constant.

The second structure is axisymmetry with swirl satisfying a quantitative bound.
The available estimates for the angular component are substantially weaker than
in the no-swirl case, and a full Liouville theorem in this setting should be
recorded as a separate hypothesis.  We formulate the hypothesis with an explicit
set of admissible exponents; no exponent range is concealed in the notation.

The third structure is uniform spatial decay in self-similar variables.  This
condition rules out spatial escape in weighted \(L^2\), but it does not by
itself furnish a complete time-dependent Liouville theorem.  We therefore state
the exact weighted Liouville hypothesis needed for the Type I exclusion and
prove the reduction from that hypothesis.

\section{Ancient Type I solutions and self-similar variables}

\begin{definition}\label{def:type-i-ancient}
A smooth solution \(U\) of \eqref{eq:ns} on
\(\mathbb R^3\times(-\infty,0)\) is called a Type I ancient solution if
\eqref{eq:type-i} holds.  Its self-similar representative is
\begin{equation}\label{eq:self-similar}
V(y,\tau)=\sqrt{-t}\,U(x,t),\qquad
Q(y,\tau)=(-t)P(x,t),\qquad
y=\frac{x}{\sqrt{-t}},\qquad
\tau=-\log(-t).
\end{equation}
\end{definition}

\begin{lemma}[Self-similar equation]\label{lem:self-similar-equation}
Let \((U,P)\) be a smooth solution of \eqref{eq:ns}, and define \((V,Q)\) by
\eqref{eq:self-similar}.  Then \(V\) is defined on
\(\mathbb R^3\times\mathbb R\) and satisfies
\begin{equation}\label{eq:leray}
\partial_\tau V-\Delta V+\frac12 y\cdot\nabla V+\frac12 V
 +V\cdot\nabla V+\nabla Q=0,\qquad \nabla\cdot V=0.
\end{equation}
Conversely, every smooth solution \((V,Q)\) of \eqref{eq:leray} determines a
smooth solution of \eqref{eq:ns} by the inverse change of variables.
\end{lemma}

\begin{proof}
Write
\[
\lambda(t)=(-t)^{-1/2},\qquad y=\lambda(t)x,\qquad
\tau=-\log(-t),
\]
so that \(U(x,t)=\lambda(t)V(y,\tau)\) and
\(P(x,t)=\lambda(t)^2Q(y,\tau)\).  Since
\[
\lambda'(t)=\frac12\lambda(t)^3,\qquad
\tau'(t)=\lambda(t)^2,\qquad
\partial_t y=\frac12\lambda(t)^2y,
\]
one has
\[
\partial_t U
=\lambda^3\left(\partial_\tau V+\frac12 y\cdot\nabla V+\frac12 V\right),
\qquad
\Delta_x U=\lambda^3\Delta_y V,
\]
and
\[
U\cdot\nabla_x U=\lambda^3V\cdot\nabla_y V,\qquad
\nabla_x P=\lambda^3\nabla_y Q.
\]
Dividing the Navier--Stokes equation by \(\lambda^3\) gives
\eqref{eq:leray}.  The divergence condition transforms as
\(\nabla_x\cdot U=\lambda^2\nabla_y\cdot V\).  The inverse computation is the
same calculation read backwards.
\end{proof}

The Type I bound is exactly the uniform \(L^\infty\) bound for \(V\):
\begin{equation}\label{eq:type-i-v}
\sup_{\tau\in\mathbb R}\|V(\cdot,\tau)\|_{L^\infty(\mathbb R^3)}<\infty .
\end{equation}
Indeed, if \(t=-e^{-\tau}\), then
\[
\|V(\cdot,\tau)\|_{L^\infty(\mathbb R^3)}
=\sqrt{-t}\,\|U(\cdot,t)\|_{L^\infty(\mathbb R^3)}.
\]

For axisymmetric flows we write
\[
x=(r\cos\theta,r\sin\theta,z),\qquad
U=U_r e_r+U_\theta e_\theta+U_z e_z.
\]
The flow is said to have no swirl if \(U_\theta\equiv0\).  The same convention is
used for \(V\), with cylindrical coordinates
\[
y=(\rho\cos\theta,\rho\sin\theta,Z).
\]

\begin{lemma}\label{lem:structure-preserved}
Axisymmetry and the no-swirl condition are preserved by the self-similar change
of variables \eqref{eq:self-similar}.  The bounds
\begin{equation}\label{eq:swirl-control}
\sup_{\tau\in\mathbb R}\sup_{\rho>0,\ Z\in\mathbb R}
\rho^\gamma |V_\theta(\rho,Z,\tau)|<\infty
\end{equation}
are invariant under self-similar time translation.
Moreover, if \(U\) satisfies \eqref{eq:type-i}, then for every \(T<0\) the
restriction of \(U\) to \(\mathbb R^3\times(-\infty,T]\) is bounded.
\end{lemma}

\begin{proof}
The self-similar change of variables is a scalar dilation in space together with
a scalar multiplication of the velocity.  It therefore commutes with rotations
around the \(z\)-axis and preserves the cylindrical decomposition of the vector
field.  This proves the assertions on axisymmetry and no swirl.  The expression
\(\rho^\gamma |V_\theta|\) is written directly in self-similar coordinates, so
\eqref{eq:swirl-control} is unchanged when \(V(y,\tau)\) is replaced by
\(V(y,\tau+a)\), \(a\in\mathbb R\).  Finally, if \(t\le T<0\), then
\eqref{eq:type-i} gives
\[
\|U(\cdot,t)\|_{L^\infty(\mathbb R^3)}
\le
\frac{1}{\sqrt{-T}}
\sup_{s<0}\sqrt{-s}\,\|U(\cdot,s)\|_{L^\infty(\mathbb R^3)} .
\]
\end{proof}

\section{The axisymmetric no-swirl case}

We use the following classical Liouville theorem in the precise form needed
below.

\begin{theorem}[Koch--Nadirashvili--Seregin--Sverak]\label{thm:knss}
Let \(W\) be a bounded smooth ancient solution of the three-dimensional
Navier--Stokes equations on \(\mathbb R^3\times(-\infty,0)\).  If \(W\) is
axisymmetric without swirl, then \(W\) is constant in space and time.
\end{theorem}

\begin{proof}
This is the axisymmetric no-swirl Liouville theorem proved in
\cite{KNSS2009}, specialized to smooth solutions.
\end{proof}

\begin{proposition}\label{prop:no-swirl-type-i}
Let \(U\) be a smooth Type I ancient solution of \eqref{eq:ns}.  If \(U\) is
axisymmetric without swirl, then \(U\equiv0\).
\end{proposition}

\begin{proof}
Let
\[
M=\sup_{t<0}\sqrt{-t}\,\|U(\cdot,t)\|_{L^\infty(\mathbb R^3)}.
\]
Fix \(T<0\), and define
\[
W_T(x,s)=U(x,T+s),\qquad s<0.
\]
Then \(W_T\) is a smooth ancient solution, axisymmetric without swirl, and
\[
\|W_T(\cdot,s)\|_{L^\infty(\mathbb R^3)}
\le \frac{M}{\sqrt{-(T+s)}}\le \frac{M}{\sqrt{-T}},
\qquad s<0.
\]
By \Cref{thm:knss}, \(W_T\) is constant on
\(\mathbb R^3\times(-\infty,0)\).  Hence \(U\) is constant on
\(\mathbb R^3\times(-\infty,T)\).  Since \(T<0\) was arbitrary and these
intervals overlap, there is a vector \(c\in\mathbb R^3\) such that
\(U(x,t)=c\) for all \(x\in\mathbb R^3\) and \(t<0\).  The Type I bound gives
\[
|c|\le \frac{M}{\sqrt{-t}}\qquad\text{for every }t<0.
\]
Letting \(t\to-\infty\) yields \(c=0\).
\end{proof}

\begin{corollary}\label{cor:no-swirl-exclusion}
No nonzero Type I ancient limit arising from a finite-time Type I singularity can
be axisymmetric without swirl.
\end{corollary}

\begin{proof}
Such a limit satisfies \eqref{eq:type-i}.  If it were axisymmetric without
swirl, \Cref{prop:no-swirl-type-i} would give \(U\equiv0\), contradicting the
nonzero normalization of the blow-up limit.
\end{proof}

\section{Axisymmetric solutions with controlled swirl}

The angular component satisfies a drift-diffusion equation with singular
coefficients on the axis.  This prevents the preceding no-swirl argument from
closing once \(U_\theta\) is nonzero.  The following formulation records exactly
the Liouville statement needed for the Type I reduction.

\begin{hypothesis}[Controlled-swirl Liouville criterion]\label{hyp:swirl}
Let \(I\) be a specified subset of \((0,\infty)\).  For every
\(\gamma\in I\), the following assertion holds.

If \(U\) is a smooth Type I ancient solution of \eqref{eq:ns}, \(U\) is
axisymmetric, and its self-similar representative \(V\) satisfies
\begin{equation}\label{eq:controlled-swirl-assumption}
\sup_{\tau\in\mathbb R}\sup_{\rho>0,\ Z\in\mathbb R}
\rho^\gamma |V_\theta(\rho,Z,\tau)|<\infty,
\end{equation}
then \(U\equiv0\).
\end{hypothesis}

\begin{remark}\label{rem:admissible-exponents}
The admissible exponent range is the set \(I\).  Thus a result proving
\Cref{hyp:swirl} with \(I=[\gamma_0,\infty)\), \(I=(0,\gamma_1]\), or any other
explicit subset immediately gives the corresponding Type I exclusion.  The
present paper does not assert an exponent range beyond the one supplied by the
hypothesis.
\end{remark}

\begin{proposition}\label{prop:swirl-exclusion}
Assume \Cref{hyp:swirl} for a specified set \(I\subset(0,\infty)\).  Let \(U\) be
a nonzero Type I ancient limit arising from a finite-time Type I singularity.  If
\(U\) is axisymmetric and its self-similar representative satisfies
\eqref{eq:controlled-swirl-assumption} for some \(\gamma\in I\), then this
configuration is impossible.
\end{proposition}

\begin{proof}
The hypotheses place \(U\) in the class covered by \Cref{hyp:swirl}.  Hence
\(U\equiv0\), contrary to the nonzero normalization of the ancient limit.
\end{proof}

\begin{remark}\label{rem:swirl-scale}
In physical variables, \eqref{eq:controlled-swirl-assumption} is equivalent to
\[
\sup_{t<0}\sup_{r>0,\ z\in\mathbb R}
(-t)^{(1-\gamma)/2}r^\gamma |U_\theta(r,z,t)|<\infty .
\]
This is the scale-invariant form of the condition.  It is written in
self-similar variables in \Cref{hyp:swirl} to avoid mixing the geometric
cylindrical variables with the Type I time weight.
\end{remark}

\section{Fast decay in self-similar variables}

We next formulate the spatial decay case.  Let \(V\) be the self-similar
representative of a Type I ancient solution.  For \(m>0\) consider the weighted
condition
\begin{equation}\label{eq:weighted-decay}
\sup_{\tau\in\mathbb R}
\int_{\mathbb R^3}(1+|y|^2)^m |V(y,\tau)|^2\,dy<\infty .
\end{equation}
This condition is stronger than local \(L^2\) boundedness; it rules out spatial
escape of \(L^2\) mass in self-similar variables.  Any integration by parts or
backward-uniqueness argument using this condition is part of the Liouville
criterion stated below.

\begin{hypothesis}[Weighted Liouville criterion]\label{hyp:weighted}
Let \(J\) be a specified subset of \((0,\infty)\).  For every \(m\in J\), the
following assertion holds.

If \(U\) is a smooth Type I ancient solution of \eqref{eq:ns} and its
self-similar representative \(V\) satisfies \eqref{eq:type-i-v} and
\eqref{eq:weighted-decay}, then \(U\equiv0\).
\end{hypothesis}

\begin{remark}\label{rem:weighted-range}
The admissible decay range is the set \(J\).  A proof of
\Cref{hyp:weighted} for a concrete interval of weights gives exactly that
interval in the Type I exclusion below.  The statement is intentionally phrased
so that the weight threshold cannot be hidden behind qualitative terminology.
\end{remark}

\begin{proposition}\label{prop:weighted-exclusion}
Assume \Cref{hyp:weighted} for a specified set \(J\subset(0,\infty)\).  Let \(U\)
be a nonzero Type I ancient limit arising from a finite-time Type I singularity.
If its self-similar representative satisfies \eqref{eq:weighted-decay} for some
\(m\in J\), then this configuration is impossible.
\end{proposition}

\begin{proof}
The hypotheses put \(U\) in the class covered by \Cref{hyp:weighted}.  Hence
\(U\equiv0\), contradicting the nonzero normalization of the ancient limit.
\end{proof}

\section{The structured Type I exclusion}

The preceding sections assemble into a single exclusion statement for structured
ancient limits.  The statement is deliberately modular: the no-swirl alternative
is unconditional, while the swirl and decay alternatives are derived from the
corresponding Liouville hypotheses with their explicit admissible parameter sets.

\begin{theorem}\label{thm:structured-exclusion}
Let \(U\) be a nonzero Type I ancient limit arising from a finite-time Type I
singularity of the three-dimensional Navier--Stokes equations, and let \(V\) be
its self-similar representative.  Then none of the following alternatives can
hold:
\begin{enumerate}
\item \(U\) is axisymmetric without swirl;
\item \Cref{hyp:swirl} holds for a specified set \(I\subset(0,\infty)\), \(U\) is
axisymmetric, and \(V\) satisfies \eqref{eq:controlled-swirl-assumption} for
some \(\gamma\in I\);
\item \Cref{hyp:weighted} holds for a specified set \(J\subset(0,\infty)\), and
 \(V\) satisfies \eqref{eq:weighted-decay} for some \(m\in J\).
\end{enumerate}
\end{theorem}

\begin{proof}
The first alternative is excluded by \Cref{cor:no-swirl-exclusion}.  The second
is excluded by \Cref{prop:swirl-exclusion}.  The third is excluded by
\Cref{prop:weighted-exclusion}.
\end{proof}

\begin{remark}\label{rem:relation-to-previous}
This theorem does not replace the compact ancient dynamics criteria proved in
the previous papers.  It supplies separate structured routes to the same
contradiction.  In the no-swirl case the route is complete.  In the swirl and
weighted cases, the remaining work is exactly the proof of the corresponding
Liouville hypothesis with an explicit admissible range.
\end{remark}

\begin{thebibliography}{9}

\bibitem{KNSS2009}
G. Koch, N. Nadirashvili, G. Seregin, and V. Sverak,
\emph{Liouville theorems for the Navier-Stokes equations and applications},
Acta Math. \textbf{203} (2009), no. 1, 83--105.

\end{thebibliography}

\end{document}
